\documentclass[12pt,a4paper]{article}

\usepackage[utf8]{inputenc}
\usepackage[T1]{fontenc}
\usepackage[brazil]{babel}
\usepackage{lmodern}
\usepackage{microtype}
\usepackage{geometry}
\usepackage{booktabs}
\usepackage{amsmath}
\usepackage{siunitx}
\usepackage{hyperref}

\geometry{margin=2.5cm}
\sisetup{locale=FR, per-mode=symbol}
\hypersetup{colorlinks=true, linkcolor=blue, urlcolor=blue}

\title{NT-001\\Preparação e Avaliação de Modelos de Previsão de Vazões}
\author{Felipe Emanuel Domiciano Ribeiro}
\date{Revisão 0 — 20 de setembro de 2026}

\begin{document}

\maketitle

\begin{center}
\begin{tabular}{ll}
\toprule
\textbf{Projeto} & Previsão diária de vazões \tabularnewline
\textbf{Fonte principal} & ANA/HIDROWEB \tabularnewline
\textbf{Horizonte inicial} & 1 dia \tabularnewline
\textbf{Unidade da vazão} & \si{\cubic\metre\per\second} \tabularnewline
\bottomrule
\end{tabular}
\end{center}

\section{Objetivo}

Esta nota técnica estabelece o procedimento mínimo para aquisição, preparação,
modelagem e avaliação de séries de vazão. O processo deve preservar a
proveniência dos dados, impedir vazamento temporal e produzir resultados
reproduzíveis.

\section{Fontes e proveniência}

Os dados hidrológicos devem ser obtidos em canais oficiais da Agência Nacional
de Águas e Saneamento Básico, preferencialmente pelo
\href{https://www.snirh.gov.br/hidroweb/}{HIDROWEB} ou pelo
\href{https://www.ana.gov.br/hidrowebservice/swagger-ui/index.html}{Hidro Webservice}.

Para cada conjunto de dados, devem ser registrados:

\begin{itemize}
    \item código e nome da estação;
    \item curso d'água, município e bacia, quando disponíveis;
    \item data de obtenção e período solicitado;
    \item tipo de série, frequência e unidade;
    \item endereço ou método de acesso;
    \item nome do arquivo original e resumo criptográfico SHA-256.
\end{itemize}

O arquivo obtido deve permanecer inalterado em
\texttt{01\_dados/input/}. Conversões e tratamentos devem gerar novos arquivos
nos respectivos diretórios de saída.

\section{Controle de qualidade}

A inspeção inicial deve identificar datas duplicadas, lacunas, valores não
numéricos, vazões negativas, qualificadores da fonte e alterações relevantes na
disponibilidade da série. Nenhuma observação pode ser excluída ou interpolada
sem justificativa registrada.

Considere a série observada

\begin{equation}
    \mathcal{Q} = \{Q_t\}_{t=1}^{N},
\end{equation}

em que $Q_t$ representa a vazão média diária, expressa em
\si{\cubic\metre\per\second}.

\section{Preparação temporal}

A divisão entre treino, validação e teste deve respeitar a ordem cronológica.
Como configuração inicial, adotam-se as proporções apresentadas na
Tabela~\ref{tab:divisao}.

\begin{table}[ht]
    \centering
    \caption{Divisão temporal inicial da série.}
    \label{tab:divisao}
    \begin{tabular}{lc}
        \toprule
        \textbf{Conjunto} & \textbf{Proporção} \\
        \midrule
        Treino & 70\% \\
        Validação & 15\% \\
        Teste & 15\% \\
        \bottomrule
    \end{tabular}
\end{table}

Os parâmetros de normalização devem ser calculados exclusivamente com o
conjunto de treino. Para média $\mu_{\mathrm{treino}}$ e desvio-padrão
$\sigma_{\mathrm{treino}}$, utiliza-se

\begin{equation}
    Q_t^{*} = \frac{Q_t - \mu_{\mathrm{treino}}}
    {\sigma_{\mathrm{treino}}}.
\end{equation}

\section{Modelo e referência}

O modelo candidato é uma rede LSTM implementada em PyTorch. A entrada inicial
corresponde aos 30 dias anteriores, e o alvo é a vazão do dia seguinte.

O desempenho deve ser comparado com o modelo de persistência:

\begin{equation}
    \widehat{Q}_{t+1} = Q_t.
\end{equation}

Alterações de arquitetura, janela, horizonte ou variáveis explicativas devem
ser registradas como decisões de modelagem.

\section{Métricas de avaliação}

As previsões devem ser avaliadas no mesmo período de teste por meio de RMSE,
MAE e eficiência de Nash--Sutcliffe (NSE):

\begin{align}
    \mathrm{RMSE} &= \sqrt{\frac{1}{n}\sum_{i=1}^{n}
    \left(Q_i-\widehat{Q}_i\right)^2}, \\
    \mathrm{MAE} &= \frac{1}{n}\sum_{i=1}^{n}
    \left|Q_i-\widehat{Q}_i\right|, \\
    \mathrm{NSE} &= 1 -
    \frac{\sum_{i=1}^{n}\left(Q_i-\widehat{Q}_i\right)^2}
    {\sum_{i=1}^{n}\left(Q_i-\overline{Q}\right)^2}.
\end{align}

O relatório deve informar o intervalo avaliado, o número de observações e as
unidades das métricas dimensionais.

\section{Reprodutibilidade e aprovação}

Cada execução oficial deve registrar código, configuração, sementes, versões
das dependências, arquivos de entrada, pesos e métricas. O relatório final deve
apresentar limitações, domínio de uso e comportamento em eventos extremos.
Resultado estatístico favorável, isoladamente, não autoriza uso operacional.

\section{Itens pendentes}

\begin{itemize}
    \item confirmar a estação fluviométrica e a bacia;
    \item definir o período útil após a inspeção da série;
    \item aprovar a política de tratamento de lacunas;
    \item definir critérios específicos para eventos extremos.
\end{itemize}

\end{document}
